\section{Positive payment and the remaining low-load source}
\label{sec:v112-source}
We now apply the joint-load theorem to the original localized
persistent source. All quantities $\varepsilon(B),h_B,r_B,\gamma_B,K_B$
are fixed at the outer stage. A load level used below classifies an
original trajectory at finite collision count; it is not a spectral,
roof, residence or comparison window.

Fix $\eta>0$ and set
\begin{equation}\label{eq:v112-paid-set}
 u_m=(2+\eta)\log(m+1),\qquad
 \mathcal Q_{B,m}=\{\mathscr L_{\varepsilon(B)}\ge u_m\}.
\end{equation}
Keep the previously paid set $\mathcal Z_{B,m}$ from
\eqref{eq:v111-collision-depth}. Define two restrictions of the
same positive source, with no new renormalization:
\begin{equation}\label{eq:v112-new-sources}
 \begin{split}
 \mathsf U_B^{\rm load}
  &=\mathsf U_B^{\rm per}[\mathcal Z_{B,m}^{\,c}\cap\mathcal Q_{B,m}],\\
 \mathsf U_B^{\rm low}
  &=\mathsf U_B^{\rm per}[\mathcal Z_{B,m}^{\,c}\cap\mathcal Q_{B,m}^{\,c}],\\
 \mathsf R_B^{\rm low}
  &=(\mathsf U_B^{\rm low}-K_BP_{r_B}^{\rm ng})_+.
 \end{split}
\end{equation}
The comparison pool is the full original nongrazing pool, not a
source-selected pool. The notation $\mathrm{low}$ refers only to
incidence load; clearance witnesses remain at arbitrary depth.

\begin{theorem}[A new collision-normalized positive payment]
\label{thm:v112-payment}
At every finite count,
\begin{equation}\label{eq:v112-positive-identity}
 \mathsf U_B^{\rm per}
 =\mathsf U_B^{\rm exc}+\mathsf U_B^{\rm load}+\mathsf U_B^{\rm low},
 \qquad
 0\le\mathsf R_B^{\rm loc}-\mathsf R_B^{\rm low}
                     \le\mathsf U_B^{\rm load}.
\end{equation}
For every fixed $B,M$, the load term has zero ordered height:
\begin{equation}\label{eq:v112-zero-height}
                \mathcal H_M(\mathsf U_B^{\rm load})=0.
\end{equation}
More explicitly, whenever $u_m>b_{\varepsilon(B)}$,
\begin{equation}\label{eq:v112-rate}
 \sup_{R,n,k}m^2\|\mathsf U_B^{\rm load}\|_\infty
 \le C_\star\exp\left\{-\eta\log(m+1)
       +2\sqrt{b_{\varepsilon(B)}(2+\eta)\log(m+1)}\right\}.
\end{equation}
Consequently
\begin{equation}\label{eq:v112-scalar}
 \mathcal H_M(\mathsf R_B^{\rm loc}-\mathsf R_B^{\rm low})=0,
 \qquad
 \mathcal E_M\le\mathcal H_M(\mathsf R_B^{\rm low})+C_MB^{-1/384}.
\end{equation}
The original assertion \eqref{eq:v76-target} is equivalent to
$\mathcal H_M(\mathsf R_B^{\rm low})\to0$ for every fixed $M$.
\end{theorem}
\begin{proof}
The three sets in \eqref{eq:v112-positive-identity} disjointly
partition the original persistent source. Thus its first equality
holds before pushforward and remains exact afterwards. Monotonicity
and the unit Lipschitz bound of the positive-part map prove the
residual inequality with the same capacity on both sides.

Lemma~\ref{lem:v111-domination} and
Theorem~\ref{thm:v112-load} bound $\mathsf U_B^{\rm load}$ by
$C_\star\exp[-u_m+2\sqrt{b_{\varepsilon(B)}u_m}]$ once
$u_m>b_{\varepsilon(B)}$. Multiplication by $m^2$, followed by
$m^2\le(m+1)^2$, gives \eqref{eq:v112-rate}. At fixed $B$ its
exponent is $-\eta\log(m+1)+O_B(\sqrt{\log(m+1)})$, tending to
minus infinity. The finitely many smaller counts do not affect the
collision limsup. This proves \eqref{eq:v112-zero-height} uniformly
in all the stated inner suprema.

Apply the residual inequality and
Corollary~\ref{cor:v111-residual} to obtain
\eqref{eq:v112-scalar}. For necessity use
$0\le\mathsf R_B^{\rm low}\le\mathsf U_B^{\rm micro}$ and the
original scalar equivalence. For sufficiency use
\eqref{eq:v112-scalar} and then send $B$ to infinity. This is the
original order of limits. The optimizing parameter in the proof of
\eqref{eq:v112-load-tail} was chosen inside a finite-count
inequality valid for every $0<\lambda<1$; it is not inserted into
a theorem which holds only at a previously fixed outer width.
\end{proof}

\begin{corollary}[Explicit support restrictions after payment]
\label{cor:v112-support}
At any regular target point contributing to $\mathsf U_B^{\rm low}$,
put $J=J_{\varepsilon(B)}$ and $N=|J|$. Then
\begin{equation}\label{eq:v112-support}
 \alpha\sum_{i\in J}d_i<u_m,\qquad
 \alpha f(N)<u_m,\qquad
 \prod_{i\in J}\frac{c_i}{2\varepsilon(B)}>e^{-u_m}.
\end{equation}
In particular $N<2\sqrt{u_m/\alpha}$ and every individual
witness has $c_i>2\varepsilon(B)e^{-u_m}$.
All these restrictions are in addition to the retained v111
incidence localization. They impose no clearance localization.
\end{corollary}
\begin{proof}
The load is strictly less than $u_m$ on the retained set. Its
nonnegative decomposition gives the first inequality. The sum of
the depths of any $N$ indices is at least $f(N)$, by the construction
in Lemma~\ref{lem:v112-depth-poly}; this proves the second. Exponentiating
minus the definition of the load proves the product inequality.
Since $f(N)\ge N^2/4$ and all load summands are nonnegative, the
remaining assertions follow, including the empty-witness case.
\end{proof}

For sufficiently large $B$, the original section cap excludes
incidence witnesses at the two endpoint collisions. The persistent
guard then still forces an interior incidence witness or an original
third-disk clearance witness. A state with $N=0$ is not removed by
this revision: the guard may be witnessed entirely by clearance.
Critical entrances and section cuts carrying the original positive
coefficient are likewise not deleted by their names.

The exact remaining target is
\begin{equation}\label{eq:v112-remaining}
 \lim_{B\to\infty}\limsup_{m\to\infty}\sup_{R,n,k}m^2
 \esssup_{t\in\mathcal C_{m,R,n,k}(M)}
       (\mathsf U_{B,m}^{\rm low}(t)-K_BP_{r_B,m}^{\rm ng}(t))_+=0.
\end{equation}
Theorems~\ref{thm:v112-load}--\ref{thm:v112-payment} prove new
finite-count concentration and positive payments in the required
normalized topology. They do not prove \eqref{eq:v112-remaining}.
Low-load shallow incidences, arbitrary-depth clearance (including
unvisited tangency), and capacity-deficient positive source remain
in that equation. The original physical section, guard, arithmetic
reference and exact-label pointwise assertion have not been weakened.
The distinction between this remaining assertion and the proved
payments is part of the theorem statement, not an inference from
software diagnostics or an editorial change of topic.
